\par\noindent\SourceLeaf{65}%
quo per $0$ diviso restat
\[
\frac{1}{2^{2h}}(-1)^{2h-1}\frac{1.2.\,.\,.\,.(2h-1)}{1.2.\,.\,.\,. 2h}\,{}^{2h}\mathfrak{B}^{(h)}\,n^{2h} = (-1)^{2h-1}\frac{1}{2^{2h}} . \frac{1}{2h}\,{}^{2h}\mathfrak{B}^{(h)}\,n^{2h}
\]
\[
= \frac{1}{2^{2h}}(-1)^{2h-1} . \frac{(2h-1)(2h-2).\,.\,.(h+1)}{1.2.\,.\,. h}\,n^{2h}
\]
\[
= \frac{1}{2^{h}}(-1)^{2h-1}\frac{1.3.5.\,.\,.(2h-1)}{1.2.3.\,.\,.\,. h.}\,\frac{n^{2h}}{2h}
\]
sive cum $(-1)^{2h-1}$, quidquid fuerit $h$, sit negativum
\[
= -\frac{1}{2^{h}}\frac{1.3.5.\,.\,.(2h-1)}{1.2.3.\,.\,. h}\,\frac{n^{2h}}{2h}.
\]
Differentiata hac expressione secundum $n$, provenit, abiecto $\mathrm{d}n$,
\[
-\frac{1}{2^{h}}\frac{1.3.5.\,.(2h-1)}{1.2.3.\,.\,. h}\,n^{2h-1},
\]
quae est manifesto quantitatis
\[
\frac{-(1-nn)^{-\frac{1}{2}}}{n}
\]
terminus $h^{\mathrm{tus}}$. Quare, cum hic non adsit, qui in illa terminus initialis
$\dfrac{-1}{n}$, illa mulctata formulae nostrae aequalis fiet, adeoque si ponatur ipsius
$\log(1+n\cos\varphi)$ coefficiens initialis $=C$ habebitur
\[
\frac{n\,{}^{*}\!\mathrm{d}C}{\mathrm{d}n} = -(1-nn)^{-\frac{1}{2}}+1.
\]
Qua aequatione, sic expressa
\[
{}^{n}\mathrm{d}C = \frac{\mathrm{d}n}{n}\bigl(1-(1-nn)^{-\frac{1}{2}}\bigr),
\]
fit integrando
\[
C = \log n - \log\frac{1-(1-nn)^{\frac{1}{2}}}{n} + \mathrm{const};
\]
est vero, cum evanescere debeat posito $n=0$, expressio const. $=\log\tfrac{1}{2}$. Quare
\[
C = \log\left(\frac{nn}{(1-\sqrt{1-nn})2}\right),
\]
aut, si malis,
\[
C = -\log\left(\frac{2(1-\sqrt{1-nn})}{nn}\right).
\]

